\section{Introduction}
\label{sec:intro}

A ground robot often faces two routes that are equally clear of hard obstacles but not equally safe to cross, with firm ground on one side and wet soil on the other. Both are geometrically feasible, yet they differ in \emph{material risk}, the soft, non-collision cost of traversing a surface or semantic class such as mud, water, or a loose shoulder. A geometry-only navigator, which reasons only about clearance to hard obstacles, treats the two routes as identical and has no channel through which material information can act on its motion. Making navigation material-aware means letting this soft cost influence how the robot moves, not only how an offline cost map is drawn.

One naive fix is to add a material-risk gradient that pushes the robot away from high-cost terrain. But a lower-risk \emph{direction} is not always a feasible \emph{maneuver}. The low-cost side may be walled off by a \emph{hard hazard}, a non-traversable or collision region that we represent here by a signed-distance field. It may also open only later, as when a blocking vehicle clears a corridor several steps after the robot first sees it. An always-active material force reacts to these apparent escapes and steers toward motions the robot cannot take, producing detours toward blocked ground or premature swerves before an escape is available. The governing question is therefore not how heavily to weight material risk, but \emph{when material information should be admitted to reshape the motion at all}. The material response should be exposed when a locally feasible lower-risk motion exists and withheld otherwise. We measure precisely this admit/withhold behavior, and it is the property we call \emph{selective}.

We instantiate the idea in a port-Hamiltonian navigation field, building on a geometry-only navigator~\citep{ellendula2025grlsnam} whose scaffold (goal attraction, hard-obstacle repulsion, and dissipation) we inherit unchanged. We augment this scaffold with two additional elements. First, a \emph{material-risk force} $f_{\mathrm{mat}}=-\nabla\tilde r$ that supplies the direction in which the soft risk $\tilde r$ acts, with a magnitude $\lambda_s(c)$ predicted from the local \emph{context} $c$ (the patch of material risk, hard-hazard distance, traversability, and goal direction the controller already sees). Second, an \emph{activation witness}, which is a local feasibility test over a few sampled short-horizon motion primitives that admits the material force only when some primitive is traversable, clears the hazard field, and improves risk by a margin. The witness is evidence that a feasible lower-risk motion exists nearby, and it is not a path the controller tracks. The field rather than the primitive generates the executed motion, and Sec.~\ref{sec:exp_rellis_dyn} reports how far the two diverge. The hard-hazard force remains active at every step; only the soft material-risk force is gated. Writing $g(c)\in[0,1]$ for the witness gate, the per-step force is
\begin{equation}
F(q)\;=\;F_{\mathrm{geom}}(q)
\;+\;\underbrace{g(c)}_{\text{exposure}}\;
     \underbrace{\lambda_s(c)}_{\text{strength}}\;
     \underbrace{f_{\mathrm{mat}}(q)}_{\text{direction}}
\;+\;\lambda_h(c)\,f_{\mathrm{haz}}(q).
\label{eq:intro_decomp}
\end{equation}
Direction, strength, and exposure are three independent factors, and keeping them separate is the point of the construction. We train the strength $\lambda_s$ with a CVaR tail-risk objective~\citep{rockafellar2000optimization,chow2015risk} so that rare, consequential transitions rather than average-case cost shape it. CVaR is a training rule for that one coefficient, not a fourth mechanism alongside direction, strength, and exposure, and we do not claim it as a source of the selectivity we report. On our head-to-head benchmark, retraining the identical architecture and gate under expected cost instead changes no outcome metric measurably (\S\ref{sec:exp_adapt}). The selective behavior comes from the gate.

% \definecolor{riskred}{RGB}{214,69,65}
% \definecolor{geomgray}{RGB}{110,110,110}
% \definecolor{execblue}{RGB}{31,78,161}
% \definecolor{matblue}{RGB}{43,111,196}
% \definecolor{passgreen}{RGB}{34,120,52}
% \definecolor{obsgray}{RGB}{120,120,120}

% \begin{figure*}[t]
% \centering
% \resizebox{\textwidth}{!}{%
% \begin{tikzpicture}[
%   font=\small,
%   obs/.style={fill=obsgray!45,draw=obsgray!80,rounded corners=1.5pt},
%   risk/.style={fill=riskred,fill opacity=0.16,draw=riskred!60,draw opacity=0.4,dashed,rounded corners=3pt},
%   start/.style={circle,fill=black,inner sep=1.6pt},
%   goal/.style={circle,draw=black,line width=0.8pt,inner sep=1.6pt},
%   bot/.style={circle,fill=execblue,inner sep=1.6pt},
%   fgeom/.style={-{Stealth[length=2mm]},line width=1.1pt,geomgray},
%   fmat/.style={-{Stealth[length=2mm]},line width=1.1pt,matblue},
%   fmatoff/.style={-{Stealth[length=2mm]},line width=1.1pt,geomgray!70,dashed},
%   fexec/.style={-{Stealth[length=2.4mm]},line width=1.7pt,execblue},
%   prim/.style={line width=0.8pt,passgreen,densely dashed},
%   primoff/.style={line width=0.7pt,gray!65,densely dotted},
%   panel/.style={draw=black!25,rounded corners=2pt,line width=0.5pt},
%   ttl/.style={font=\footnotesize\bfseries},
%   gate/.style={draw=black!55,rounded corners=1.5pt,inner sep=2pt,font=\footnotesize},
% ]
% \def\pitch{5.5}

% % ---- Panel A ----
% \begin{scope}[shift={(0,0)}]
%   \draw[panel] (-0.15,-0.35) rectangle (4.15,3.55);
%   \node[ttl,anchor=north west] at (-0.15,4.2) {(A) Geometry misses material risk};
%   \node[obs,minimum width=0.9cm,minimum height=1.0cm] at (2.0,1.6) {};
%   \fill[risk] (1.1,2.5) rectangle (3.0,3.25);
%   \node[riskred!80!black,font=\scriptsize] at (2.05,2.88) {material risk};
%   \node[start,label={[font=\scriptsize]below:start}] (sa) at (0.35,1.6) {};
%   \node[goal,label={[font=\scriptsize]below:goal}] (ga) at (3.75,1.6) {};
%   \draw[fgeom] (sa) .. controls (1.4,3.0) and (2.6,3.0) .. (ga);
%   \draw[fgeom,geomgray!55] (sa) .. controls (1.4,0.3) and (2.6,0.3) .. (ga);
%   \node[geomgray!90,font=\scriptsize] at (0.95,2.55) {feasible};
%   \node[geomgray!90,font=\scriptsize] at (0.95,0.55) {feasible};
%   \draw[fmat] (2.05,2.35) -- ++(0,-0.7) node[right,font=\scriptsize,matblue]{$-\nabla\tilde r$};
% \end{scope}

% % ---- Panel B ----
% \begin{scope}[shift={(\pitch,0)}]
%   \draw[panel] (-0.15,-0.35) rectangle (4.15,3.55);
%   \node[ttl,anchor=north west] at (-0.15,4.2) {(B) Lower risk $\ne$ feasible};
%   \node[anchor=north west,font=\scriptsize,gray!55!black] at (-0.05,3.45) {escape blocked};
%   \node[obs,minimum width=0.9cm,minimum height=0.9cm] at (2.05,2.05) {};
%   \node[obs,minimum width=1.6cm,minimum height=0.5cm,fill=obsgray!70] at (2.25,0.45) {};
%   \node[font=\scriptsize,white] at (2.25,0.45) {blocked};
%   \node[bot,label={[font=\scriptsize]left:robot}] (rb) at (0.7,1.55) {};
%   \node[goal] (gb) at (3.8,1.6) {};
%   \draw[primoff] (rb) .. controls (1.4,0.9) and (2.0,0.75) .. (2.45,0.9);
%   \draw[primoff] (rb) .. controls (1.5,1.5) and (1.9,1.85) .. (2.15,1.98);
%   \draw[primoff] (rb) .. controls (1.4,2.2) and (2.0,2.65) .. (2.35,2.7);
%   \node[gray!65,font=\scriptsize,anchor=west] at (0.5,0.78) {candidates};
%   \draw[fgeom] (rb) -- ++(1.05,0.12) node[above,font=\scriptsize,geomgray]{$F_{\mathrm{geom}}$};
%   \draw[fmatoff] (rb) -- ++(0.18,-0.9) node[below,font=\scriptsize,gray!70,xshift=-1mm]{$\lambda_s f_{\mathrm{mat}}$};
%   \draw[fexec] (rb) -- ++(1.05,0.04);
%   \node[gate] at (2.05,3.05) {witness fails: $g=0$};
% \end{scope}

% % ---- Panel C ----
% \begin{scope}[shift={(2*\pitch,0)}]
%   \draw[panel] (-0.15,-0.35) rectangle (4.15,3.55);
%   \node[ttl,anchor=north west] at (-0.15,4.2) {(C) Witness gates material force};
%   \node[anchor=north west,font=\scriptsize,gray!55!black] at (-0.05,3.45) {escape open};
%   \node[obs,minimum width=0.9cm,minimum height=0.9cm] at (2.05,2.05) {};
%   \node[bot] (rc) at (0.7,1.55) {};
%   \node[goal] (gc) at (3.8,1.6) {};
%   \draw[primoff] (rc) .. controls (1.4,2.2) and (2.0,2.65) .. (2.35,2.7);
%   \draw[primoff] (rc) .. controls (1.5,1.5) and (1.9,1.85) .. (2.15,1.98);
%   \draw[prim] (rc) .. controls (1.5,0.5) and (2.6,0.45) .. (3.3,0.9)
%      node[pos=0.62,below=0.5pt,font=\scriptsize,passgreen]{witness $\checkmark$};
%   \draw[fgeom] (rc) -- ++(1.05,0.12) node[above,font=\scriptsize,geomgray]{$F_{\mathrm{geom}}$};
%   \draw[fmat] (rc) -- ++(0.28,-0.9) node[below,font=\scriptsize,matblue,xshift=-1mm]{$\lambda_s f_{\mathrm{mat}}$};
%   \draw[fexec] (rc) .. controls (1.3,1.0) and (2.2,1.0) .. (2.9,1.4);
%   \node[execblue,font=\scriptsize] at (2.4,1.62) {executed};
%   \node[gate] at (2.05,3.05) {witness passes: $g=1$};
% \end{scope}

% % ---- decomposition banner ----
% \node[draw=black!30,rounded corners=2pt,fill=black!3,inner sep=4pt,font=\small]
%   at ({\pitch+2.0},-1.05)
%   {$F(q)\;=\;\underbrace{F_{\mathrm{geom}}(q)}_{\text{geometry}}\;+\;\underbrace{g(q)}_{\text{exposure}}\,\underbrace{\lambda_s(c)}_{\text{strength}}\,\underbrace{f_{\mathrm{mat}}(q)}_{\text{direction}}\;+\;\underbrace{\lambda_h(c)\,f_{\mathrm{haz}}(q)}_{\text{always-on hazard}}$};
% \end{tikzpicture}}
% \caption{\textbf{Where material risk is admitted, and where it is withheld.}
% \textbf{(A)} Two routes are equally clear of hard obstacles but differ in material risk; the material-risk force points away from the high-cost region.
% \textbf{(B)} When the lower-risk side is blocked, no sampled primitive clears the hazard field, the witness fails ($g{=}0$), and the soft force is suppressed while the executed motion stays geometry-driven. \textbf{(C)} When an escape is open, one primitive passes the witness ($g{=}1$) and admits the soft force; the sampled primitive is only \emph{evidence} of feasibility: the executed trajectory (solid) is generated by the field and need not coincide with it. The hard-hazard force is active throughout.}
% \label{fig:teaser}
% \end{figure*}

Because material risk enters as its own force term rather than as one summand of an undifferentiated objective, its direction, magnitude, and exposure can each be read off and tested separately. This makes material-aware behavior falsifiable rather than merely inferred from final trajectories, and it lets us state three consequences of the construction (Prop.~\ref{prop:enrichment}). The enriched field stays close to the geometry-only field when the material force is small (C1), produces no lateral escape when none is admitted (C2), and deflects toward a feasible lower-risk direction when the witness admits one (C3). We test these across the three regimes used throughout the paper, in which a feasible lower-risk alternative exists, the lower-risk-looking direction is blocked or not yet open, or the scene is risk-neutral so that geometry-only behavior should be preserved. The test itself runs through the activation and suppression rates (CAR/FAR/SR) that the separable force makes directly observable.


We make three contributions.
\begin{itemize}
  \item \textbf{A separable, gated material-risk force.} We give material risk its own directional force channel in a navigation field, exposed by a local feasibility witness that admits it only when a traversable lower-risk motion exists nearby. The witness supplies evidence of feasibility, not a reference trajectory the controller tracks.
  \item \textbf{Context-conditioned magnitude that transfers across domains.} A learned coefficient rescales the material force from local context, adapting by roughly an order of magnitude between off-road and semantic-map settings without per-domain retuning.
  \item \textbf{Mechanistic evaluation of material influence.} We measure when the force activates and suppresses (CAR/FAR/SR, false pre-activation), how far the executed motion departs from the witness primitive, the isolated behavioral effect of the soft coefficient, and robustness to perception degradation.
\end{itemize}

We use a port-Hamiltonian field because it gives a clean additive force decomposition, which is what makes the three factors separately measurable. Nothing in the witness-gated construction depends on that choice. The gate reads only the local context and the sampled primitives, and it multiplies a single additive term, so it should transfer to other vector-field navigators such as artificial potential fields~\citep{khatib1986real}, Riemannian motion policies, or geometric fabrics. We report evidence that the mechanism transfers across \emph{domains}, namely off-road, semantic-map, dynamic, and highway settings. We do not test it in another navigator family, so its portability across field formulations remains an expectation rather than a result.